\documentclass[12pt]{article}
\usepackage[utf8]{inputenc}
\usepackage[T1]{fontenc}
\usepackage{hyperref}
\title{Watiwa Language Resources for K-12 Curriculum Development}
\author{Compiled by Claude Sonnet 5.5 (claude-sonnet-5-5) for Andr\'as Kornai}
\date{}
\begin{document}
\maketitle

\section{Introduction}
Watiwa (ISO 639-3 wtf; also called Dumpu or Watifa) is a Rai Coast language of Madang Province, Papua New Guinea. It has about 500 speakers across six villages, including Dumpu. According to Wikipedia, Tok Pisin is used day to day and the language is used largely as a secret code when outsiders are nearby. Wikipedia projects extinction within a few decades.

\section{Primary Educational Resources}
I found no primers, textbooks or school programmes.

\section{Linguistic and Reference Materials}
\begin{itemize}
\item Wikipedia, Watiwa language: \url{https://en.wikipedia.org/wiki/Watiwa_language}
\item Wikipedia, Dumpu village: \url{https://en.wikipedia.org/wiki/Dumpu,_Papua_New_Guinea}
\end{itemize}
I found no grammar or dictionary.

\section{Religious and Community Resources}
\begin{itemize}
\item Joshua Project people-group profile: \url{https://joshuaproject.net/people_groups/15889}
\item Joshua Project language page: \url{https://joshuaproject.net/languages/wtf}
\end{itemize}
I found no scripture.

\section{Digital and Audio Resources}
I found none.

\section{Curriculum Development Context}
Almost nothing is documented. A curriculum would need community-led documentation first, and would likely sit alongside Tok Pisin and English instruction.
\end{document}